% ============================================================
%  03_dac_ta_he_thong.tex  --  Dac ta he thong va luong du lieu
% ============================================================

\section{Đặc tả hệ thống và luồng dữ liệu}
\label{sec:dacta}

\subsection{Kiểu dữ liệu trung tâm}

Module Inventory xây dựng xung quanh kiểu dữ liệu \textit{Book}, được định nghĩa
giống nhau trong cả hai phiên bản:

\begin{lstlisting}[language=C, caption={Định nghĩa struct Book}]
#define TITLE_MAX  32
#define AUTHOR_MAX 32

typedef struct {
    char title[TITLE_MAX];        /* offset  0 .. 31 */
    char author[AUTHOR_MAX];      /* offset 32 .. 63 */
    int  year;                    /* offset 64 .. 67 */
    int  copies_available;        /* offset 68 .. 71 */
} Book;
\end{lstlisting}

Struct có kích thước 72 byte. Hai mảng ký tự 32 byte nằm liền kề là điểm nóng
của lỗi CWE-121: nếu ghi vào \textit{title} không giới hạn độ dài, byte thứ 32
trở đi sẽ ghi đè lên \textit{author}.

\subsection{Phiên bản vulnerable -- \textit{vulnerable\_inventory.c}}

File này cố ý chứa ba lỗi bộ nhớ để phục vụ mục đích giáo dục. Ba hàm chứa lỗi
là:

\subsubsection{(A) CWE-121 -- \textit{add\_book\_title\_UNSAFE\_HEAP()}}

\begin{lstlisting}[language=C, caption={Lỗi CWE-121: strcpy không giới hạn độ dài (heap variant)}]
char *add_book_title_UNSAFE_HEAP(const char *title) {
    char *buf = (char *)malloc(TITLE_MAX_HEAP); /* cap phat 32 byte */
    if (!buf) return NULL;
    strcpy(buf, title); /* !!! KHONG kiem tra do dai nguon !!! */
    return buf;
}
\end{lstlisting}

Nếu \textit{strlen(title) >= 32}, \textit{strcpy} ghi vượt ra ngoài 32 byte đã
cấp phát, đè lên metadata heap hoặc dữ liệu của allocation khác.

\subsubsection{(B) CWE-416 -- \textit{demo\_use\_after\_free()}}

\begin{lstlisting}[language=C, caption={Lỗi CWE-416: truy cập con trỏ sau free()}]
void demo_use_after_free(void) {
    Book *b = (Book *)malloc(sizeof(Book));
    if (!b) return;
    strcpy(b->title, "Clean Code");
    b->year = 2008;
    free(b);
    /* LOI: b->year duoc doc sau khi da free(b) */
    printf("  [UAF] nam xuat ban sau khi free: %d\n", b->year);
}
\end{lstlisting}

Sau \textit{free(b)}, vùng nhớ có thể được tái cấp phát cho mục đích khác. Đọc
\textit{b->year} là hành vi không xác định (Undefined Behavior).

\subsubsection{(C) CWE-401 -- \textit{demo\_memory\_leak()}}

\begin{lstlisting}[language=C, caption={Lỗi CWE-401: vòng lặp cấp phát thiếu free()}]
void demo_memory_leak(int count) {
    for (int i = 0; i < count; i++) {
        Book *b = (Book *)malloc(sizeof(Book));
        if (!b) continue;
        b->year = 2000 + i;
        /* LOI: thieu free(b) -- bo nho bi ro ri sau moi vong lap */
    }
}
\end{lstlisting}

Mỗi vòng lặp cấp phát một đối tượng \textit{Book} (72 byte), gán lại cho con
trỏ cục bộ \textit{b} ở vòng tiếp theo mà không giải phóng vùng nhớ cũ.
Sau 5 lần gọi: 5 $\times$ 72 = 360 byte rò rỉ.

\subsection{Phiên bản secure -- \textit{secure\_inventory.c}}

Mỗi lỗi được sửa theo nguyên tắc lập trình C an toàn:

\subsubsection{(A) Sửa CWE-121 -- dùng \textit{snprintf}}

\begin{lstlisting}[language=C, caption={Bản vá CWE-121: snprintf giới hạn độ dài}]
char *add_book_title_SAFE_HEAP(const char *title) {
    char *buf = (char *)malloc(TITLE_MAX_HEAP);
    if (!buf) return NULL;
    snprintf(buf, TITLE_MAX_HEAP, "%s", title); /* an toan: tu dong cat bot */
    return buf;
}

int add_book_title_SAFE(Book *b, const char *title) {
    if (b == NULL || title == NULL) {
        fprintf(stderr, "[LOI] Tham so NULL\n");
        return -1;
    }
    int written = snprintf(b->title, TITLE_MAX, "%s", title);
    return (written >= TITLE_MAX) ? 1 : 0; /* tra ve 1 neu bi cat bot */
}
\end{lstlisting}

\textit{snprintf} nhận tham số \textit{size} là kích thước tối đa (bao gồm
\textit{'\textbackslash0'}), đảm bảo không bao giờ ghi quá giới hạn. Nếu chuỗi
nguồn dài hơn, nó bị cắt bớt an toàn (\textit{truncated}) thay vì gây tràn.

\subsubsection{(B) Sửa CWE-416 -- sao chép trước, NULL hóa sau}

\begin{lstlisting}[language=C, caption={Bản vá CWE-416: sao chép dữ liệu trước free()}]
void demo_use_after_free_FIXED(void) {
    Book *b = (Book *)malloc(sizeof(Book));
    if (!b) { fprintf(stderr, "[LOI] malloc that bai\n"); return; }
    add_book_title_SAFE(b, "Clean Code");
    b->year = 2008;
    int year_copy = b->year;   /* SAO CHEP truoc khi free */
    free(b);
    b = NULL;                  /* TRANH dangling pointer */
    printf("  [FIXED] nam xuat ban (sao chep truoc free): %d\n", year_copy);
}
\end{lstlisting}

\subsubsection{(C) Sửa CWE-401 -- giải phóng đầy đủ}

\begin{lstlisting}[language=C, caption={Bản vá CWE-401: free() toàn bộ trong vòng lặp}]
void demo_memory_leak_FIXED(int count) {
    Book **books = (Book **)malloc(sizeof(Book *) * (size_t)count);
    if (!books) { fprintf(stderr, "[LOI] malloc that bai\n"); return; }
    for (int i = 0; i < count; i++) {
        books[i] = (Book *)malloc(sizeof(Book));
        if (!books[i]) continue;
        books[i]->year = 2000 + i;
    }
    for (int i = 0; i < count; i++) {
        free(books[i]);   /* giai phong tung phan tu */
        books[i] = NULL;
    }
    free(books);           /* giai phong mang con tro */
}
\end{lstlisting}

\subsection{Luồng dữ liệu tổng thể}

Luồng thực hiện xuyên suốt toàn bộ bài thực hành:

\begin{enumerate}
    \item Biên dịch: GCC biên dịch cả hai file C với tùy chọn
          \textit{-g -Wall -Wextra -fsanitize=address}, tạo hai binary có ASan
          được nhúng sẵn.
    \item Chạy kịch bản: Ba kịch bản (\textit{overflow}, \textit{uaf},
          \textit{leak}) được chạy lần lượt trên \textit{vulnerable\_inventory}.
          ASan hoặc LeakSanitizer ngắt chương trình và in thông báo lỗi chi tiết.
    \item Đối chiếu bản secure: Cùng ba kịch bản chạy trên
          \textit{secure\_inventory} -- kết thúc bình thường (exit code 0), không
          có thông báo lỗi.
    \item Đánh giá rủi ro: \textit{python/risk\_assessor.py} đọc
          \textit{dataset/vulnerabilities.json}, tính điểm CVSS-lite cho năm lỗ
          hổng, sắp xếp theo điểm giảm dần và ghi kết quả ra
          \textit{dataset/risk\_assessment\_results.json}.
    \item Kiểm thử tự động: \textit{tests/test\_memory\_safety.py} tổng
          hợp cả hai luồng trên, so sánh kết quả với baseline và xác nhận tính
          nhất quán.
\end{enumerate}

\begin{table}[H]
\centering
\begin{tabularx}{\textwidth}{|p{3.5cm}|X|p{3.5cm}|}
\hline
\textbf{Thành phần} & \textbf{Mô tả} & \textbf{Đầu ra} \\ \hline
\textit{vulnerable\_inventory.c} & Mã nguồn C với 3 lỗi cố ý & Binary + ASan output \\ \hline
\textit{secure\_inventory.c} & Bản đã sửa 3 lỗi & Binary sạch \\ \hline
\textit{vulnerabilities.json} & Dataset 5 lỗ hổng với thuộc tính CVSS-lite & Đầu vào đánh giá rủi ro \\ \hline
\textit{risk\_assessor.py} & Tính điểm và xếp hạng rủi ro & \textit{risk\_assessment\_results.json} \\ \hline
\textit{test\_memory\_safety.py} & 8 bài kiểm thử tự động & 8/8 PASS \\ \hline
\end{tabularx}
\caption{Các thành phần của hệ thống thực hành}
\end{table}
